\section{总结与延伸}

\subsection{本章小结}

这一讲的主结论可以压缩成四句：
\begin{enumerate}
    \item \textbf{RLHF 有用，但很难规模化}，而且容易出现 overoptimization 和 calibration 退化。
    \item \textbf{PPO 能工作，但工程上重}，value function、rollout、reward shaping 都让系统复杂。
    \item \textbf{GRPO 把 PPO 进一步简化}，去掉 value 网络，用 group normalization 代替复杂的 advantage 估计。
    \item \textbf{RLVR 是 reasoning 模型的重要路线}，DeepSeek R1、Kimi K1.5、Qwen 3 都说明了这点。
\end{enumerate}

\begin{importantbox}{今天最值得带走的一句话}
当 reward 可验证时，RL 不再只是“更难调的偏好优化”，它会变成一个能真正推动 reasoning 模型扩展的训练工具。
\end{importantbox}

\end{document}
